\section{Conclusion}
\label{sec:conclusion}

With lineage-grouped evaluation, frozen ESM3 representations show an
improvement over sequence features for CATNAP bnAb
classification and show a modest increase in breadth performance. However,
absolute breadth prediction remains weak, indicating that heavy-chain sequence
features alone, even from large PLMs, are insufficient for accurately modeling
functional neutralization breadth. Furthermore, our unsupervised OAS LoRA
adaptation run did not improve downstream performance, suggesting that
general-repertoire adaptation without target-specific supervision is
insufficient for capturing the complex somatic hypermutation patterns critical
to HIV-1 neutralization.

To our knowledge the framing itself is novel: prior machine learning for HIV
neutralization is predominantly \emph{virus-centric}, scoring envelope
resistance against a fixed antibody
panel~\cite{danaila2022review,igiraneza2024multitask,elanbari2025brave}, and the
closest antibody-centric breadth estimator infers breadth by molecular modeling
onto a template antibody--antigen complex for a single epitope
class~\cite{conti2019breadth}. We instead classify breadth directly from
heavy-chain sequence (rather than from a modeled structure), epitope-agnostically,
and under a lineage-held-out split. In
that out-of-lineage setting frozen ESM3 raises pooled AUROC from $0.625$ for a
handcrafted biophysical baseline to $0.746$ (paired $+0.121$, group-bootstrap
interval $[0.012,0.228]$ excluding zero). Because no published benchmark shares
this sequence-only, epitope-agnostic, lineage-held-out protocol, external AUROCs
from random-split or virus-side tasks are not directly comparable; the
meaningful baseline is our own handcrafted-feature model, over which ESM3 shows a
credible out-of-lineage gain.

The residue-level VAE interpolation analysis shows cautionary results. The
per-residue VAE model introduces substantial reconstruction error and fails to
outperform simpler midpoint controls on the VRC34 lineage, likely
due to the obvious lack of bnAb data. Due to that fact, we cannot draw
strong conclusions on weither a VAE would meaningfully
support the generation or design of functional antibodies, or
about VAE-space interpolation.
While the LoRA maturation correlations were higher,
the limited exploratoon of this tuning
does not establish a negative causal effect. 

Ultimately, our strongest contribution is methodological; implementing
lineage-aware splits, pooled-versus-fold validation reporting, and
group-resampled uncertainty estimation changes the interpretation
of these experiments.

Our results motivate natural next steps. Because
heavy-chain sequence alone proved insufficient for breadth, richer inputs such
as paired heavy/light chains or explicit structure are worth exploring.
Additionally,
the results of the VAE and adaptation analyses were limited primarily by the scarcity
of labeled bnAb data, larger cohorts and external validation on held-out
lineages are needed before stronger conclusions can be drawn.
